\NativeSlot{50}{88}{%
% p070_s01

\ExamQuestion{6.}{5}%
钟面上指针旋转一大格是\mbox{（\hspace{1.8em}）}度，指针从“2”到“5”是指针\mbox{（\hspace{1.8em}）}时针旋转\mbox{（\hspace{1.8em}）}度。\par

\par\noindent%
\ExamTile{.56}{%
% p070_s02

\ExamQuestion{7.}{5}%
如图，三角尺向下平移了\mbox{（\hspace{1.8em}）}毫米。\par
}\hfill%
\ExamTile{.415}{%
\BookFigure{figures/p070_s02-01.tex}
}\par

\ExamGroup{二、判断。{\normalfont\kaishu\fontsize{10.8bp}{14bp}\selectfont （每题2分，共10分）}}

\par\noindent%
\ExamTile{.485}{%
% p070_s03

\ExamQuestion{1.}{5}%
判断：三角形都有三条对称轴。\mbox{（\hspace{1.8em}）}\par
}\hfill%
\ExamTile{.485}{%
% p070_s04

\ExamQuestion{2.}{5}%
判断：物体平移时，物体上的每一个点走过的路程一样长。\mbox{（\hspace{1.8em}）}\par
}\par
}
\NativeSlot{198}{36}{%
\par\noindent%
\ExamTile{.485}{%
% p070_s05

\ExamQuestion{3.}{5}%
判断：图形经过旋转后，形状会发生改变。\mbox{（\hspace{1.8em}）}\par
}\hfill%
\ExamTile{.485}{%
% p070_s06

\ExamQuestion{4.}{5}%
判断：正方形的对角线是它的对称轴。\mbox{（\hspace{1.8em}）}\par
}\par
}
\NativeSlot{258}{36}{%
\par\noindent%
\ExamTile{.56}{%
% p070_s07

\ExamQuestion{5.}{5}%
判断：把图中的图形在对称轴的另一边补充完整，这图形是等腰梯形。\mbox{（\hspace{1.8em}）}\par
}\hfill%
\ExamTile{.415}{%
\BookFigure{figures/p070_s07-01.tex}
}\par
}
\NativeSlot{294}{92}{%
\ExamGroup{三、选择（将正确答案的序号填在括号里）。{\normalfont\kaishu\fontsize{10.8bp}{14bp}\selectfont （每题2分，共16分）}}

\par\noindent%
\ExamTile{.485}{%
% p070_s08

\ExamQuestion{1.}{5}%
下列图形中，对称轴最多的是\mbox{（\hspace{1.8em}）}，对称轴最少的是\mbox{（\hspace{1.8em}）}。\par
\BookFigure{figures/p070_s08-01.tex}
}\hfill%
\ExamTile{.485}{%
% p070_s09

\ExamQuestion{2.}{5}%
下面不是轴对称图形的是\mbox{（\hspace{1.8em}）}。\par
\ExamOptions{A. 长方形}{B. 平行四边形}{C. 圆}{D. 半圆}
}\par
}
\NativeSlot{455}{43}{%
% p070_s10

\ExamQuestion{3.}{5}%
右边的轴对称图形是从纸\mbox{（\hspace{1.8em}）}上剪下来的。选A、B或C。\par
\BookFigure{figures/p070_s10-01.tex}
}
\NativeSlot{530}{36}{%
% p070_s11

\ExamQuestion{4.}{5}%
一个三角形任意一条边上的高所在的直线都是它的对称轴，这个三角形是\mbox{（\hspace{1.8em}）}。\par
\ExamOptions{A. 等腰三角形}{B. 等腰直角三角形}{C. 等边三角形}{D. 任意三角形}
}
\NativeSlot{580}{128}{%
\par\noindent%
\ExamTile{.485}{%
% p070_s12

\ExamQuestion{5.}{5}%
小东在描述一个图形：它有四条边，没有平行线，有一条对称轴，以下\mbox{（\hspace{1.8em}）}是小东描述的。选A、B、C或D。\par
\BookFigure{figures/p070_s12-01.tex}
}\hfill%
\ExamTile{.485}{%
% p070_s13

\ExamQuestion{6.}{5}%
下面四幅图中，\mbox{（\hspace{1.8em}）}平移到位置M后能与N成轴对称图形。选A、B、C或D。\par
\BookFigure{figures/p070_s13-01.tex}
}\par
}
